\section{引言}

随着在轨服务、空间碎片清理和失效航天器维护需求的增长，空间机械臂已成为执行非合作目标接近、捕获与操作任务的重要装备。与地面固定基座机械臂不同，捕获前空间机械臂通常处于自由漂浮或自由飞行状态，基座姿态和位置会受到机械臂关节运动的反作用影响；同时，目标捕获任务对末端轨迹精度、关节运动平滑性和系统鲁棒性均提出较高要求。因此，如何在基座不直接受控的条件下实现高精度轨迹跟踪，是空间机械臂运动控制中的关键问题之一\cite{yan}。

自由漂浮空间机械臂的研究通常从系统建模、轨迹规划和闭环控制三个层面展开。已有工作针对自由漂浮模式下的运动学冗余、动力学奇异和关节速度约束问题，采用粒子群优化等方法生成满足约束的关节轨迹\cite{wang2015}；针对捕获前反作用抑制和参数不确定问题，也有研究将反作用零空间规划、时延估计和自适应滑模控制结合，以降低对精确动力学模型的依赖\cite{yu2022}。近年来，固定时间、预定义时间和指定时间控制进一步受到关注，其优势在于可在设计阶段给出与初始误差弱相关或无关的收敛时间估计。相关研究已经将预定义时间稳定思想用于自由飞行空间机器人 SE(3) 位姿跟踪\cite{xu2024}，也将固定时间容错控制用于含执行器故障和性能约束的自由漂浮空间机械臂\cite{li2025}。

不过，空间机械臂的非线性强耦合动力学使固定时间控制器设计并不直接。若直接在原非线性系统上叠加高增益或奇异型指定时间反馈，往往会带来控制输入峰值大、参数调节困难和工程实现敏感等问题。Zhou、Michiels 和 Chen 提出的光滑周期延迟反馈方法表明，对于可控线性时滞系统，利用周期延迟反馈可以实现固定时间镇定，并且反馈增益可选为连续、连续可微甚至光滑形式，从而避免传统指定时间高增益方法在终端时刻出现的增益奇异\cite{zhou2022}。该结论为固定时间控制提供了一种不同于非光滑终端滑模的设计思路。

需要强调的是，文献\cite{zhou2022}的严格结论主要面向线性可控时滞系统，不能未经处理直接套用于自由漂浮空间机械臂的原始非线性动力学。基于这一点，本文采用“动力学补偿--误差线性化--周期延迟反馈”的结构：先利用动量守恒关系将捕获前自由漂浮空间机械臂约化为等效关节空间模型，再通过计算力矩补偿得到从辅助加速度到跟踪误差的线性可控通道，最后在该误差通道中嵌入光滑周期延迟反馈。这样既保留 PDF 方法的固定时间镇定机理，又使其作用对象与原理论假设保持一致。

本文主要工作如下：第一，建立捕获前自由漂浮空间机械臂的等效关节空间动力学模型，并明确基座反作用对末端运动学的影响；第二，构造面向 PDF 设计的误差通道规范化方法，将关节轨迹跟踪问题转化为线性可控误差系统镇定问题；第三，给出结合动力学补偿和周期延迟误差反馈的关节力矩控制律，并分析理想补偿和有界扰动条件下的闭环性质；第四，基于七自由度空间机械臂参数搭建 MATLAB 仿真，验证所提控制结构的可运行性和跟踪性能。
